\documentclass[11pt]{article}

% ---- Packages ----
\usepackage[utf8]{inputenc}
\usepackage{amsmath, amssymb, amsthm}
\usepackage{mathtools}
\usepackage{bm}
\usepackage[margin=1in]{geometry}
\usepackage{hyperref}
\usepackage{booktabs}
\usepackage{array}
\usepackage{enumitem}
\usepackage{listings}
\usepackage{xcolor}

\hypersetup{
  colorlinks=true,
  linkcolor=blue!50!black,
  citecolor=blue!50!black,
  urlcolor=blue!50!black
}

% ---- Theorem environments ----
\newtheorem{proposition}{Proposition}
\newtheorem{lemma}{Lemma}
\newtheorem{corollary}{Corollary}
\theoremstyle{definition}
\newtheorem{definition}{Definition}
\newtheorem{assumption}{Assumption}
\newtheorem{remark}{Remark}

% ---- Notation shortcuts ----
\newcommand{\R}{\mathbb{R}}
\newcommand{\E}{\mathbb{E}}
\newcommand{\Var}{\operatorname{Var}}
\newcommand{\Cov}{\operatorname{Cov}}
\renewcommand{\Pr}{\mathbb{P}}
\newcommand{\Norm}{\mathcal{N}}
\newcommand{\D}{\mathcal{D}}
\newcommand{\Ltot}{\mathcal{L}}
\newcommand{\indicator}{\mathbf{1}}
\newcommand{\given}{\,\vert\,}

% ---- Listings setup for R code ----
\lstset{
  basicstyle=\ttfamily\footnotesize,
  breaklines=true,
  showstringspaces=false,
  language=R,
  commentstyle=\color{gray},
  keywordstyle=\color{blue!60!black},
  stringstyle=\color{green!50!black},
  frame=single,
  framerule=0.3pt
}

% ---- Title block ----
\title{
  AR(2) in a Spatial Skew-$t$ Model:\\
  Why a Temporal Prior Loses to an i.i.d.\ Baseline\\
  Under Spatial-Only Holdout
}
\author{Technical Note --- \texttt{ar2\_rethink}}
\date{}

\begin{document}
\maketitle

% ====================================================================
\begin{abstract}
We document a sequence of derivations explaining why an AR(2) temporal
prior on the latent processes of a spatial skew-$t$ model fails to
dominate its independent-in-time counterpart when the prediction task
is scored by the Brier score on held-out \emph{spatial} sites. The
construction normalizes each AR(2) latent so that its stationary
marginal variance equals one; we show this construction implies the
time-$t$ marginal of every latent is invariant to the autoregressive
coefficients $\phi_1, \phi_2$. The Brier loss in the spatial-holdout
design integrates the predictive kernel against precisely this
time-$t$ marginal, so the temporal coupling marginalizes out. The AR(2)
parameters can still leak into the predictive through the posterior of
the latents at training times, but each of the four available leakage
channels is damped (discrete knot membership, likelihood-dominated
latents, $\phi$-independent spatial noise), while the cost of six
additional AR coefficients inflates posterior variance. We close with
three concrete redesigns---block-time-forecast holdout, multivariate
energy score, variogram score---under which the AR(2) prior is
expected to recover its advantage.
\end{abstract}

% ====================================================================
\section{Introduction}
\label{sec:intro}

The simulation study in
\path{code/analysis/simstudy/} compares eight prediction methods for
spatial extremes, including four variants of the skew-$t$ knot model
of which two (methods 7 and 8) place an AR(2) prior on every latent
temporal process. Tabulated Brier and quantile scores for settings 9
and 10 show that the AR(2) methods do \emph{not} systematically
outperform their i.i.d.\ siblings; in setting 10 the
\verb|K=5| AR(2) variant is in fact worse than the
\verb|K=5| no-AR baseline at every quantile.

This note explains the discrepancy. The conclusion has nothing to do
with bugs in the AR(2) implementation in
\verb|simulate_ar2_standard()| and everything to do with the
interaction between (i) the AR(2) standardisation, (ii) the spatial
structure of the holdout, and (iii) the scoring rule. The reasoning
is reproduced here in a single place because it determines what
experiments would actually credit the AR(2) machinery.

The structure of the note is as follows. Section~\ref{sec:model}
fixes notation and lists the assumptions of the underlying model.
Section~\ref{sec:ar2-construction} formalises the AR(2)
standardisation and establishes the marginal invariance proposition.
Section~\ref{sec:derivations} carries the invariance through the
Brier predictive integral and identifies the four leakage channels.
Section~\ref{sec:redesign} proposes alternative scoring designs.
Section~\ref{sec:conclusion} summarises.

% ====================================================================
\section{Model Setup and Notation}
\label{sec:model}

\subsection{Spatial skew-$t$ knot model}

Let $\{s_i\}_{i=1}^{n_s} \subset \R^2$ denote spatial locations and
$t \in \{1, \dots, T\}$ time indices. Fix $K$ knots in the spatial
domain. The conditional model for the observed process is

\begin{equation}
  Y_t(s) = x_t(s)^\top \beta
         + \lambda \, z_{g(s, t), t}
         + \frac{1}{\sqrt{\tau_{g(s, t), t}}} \,
           \big[L_C \, \varepsilon_t\big]_s,
  \qquad \varepsilon_t \sim \Norm(0, I_{n_s}),
  \label{eq:kernel}
\end{equation}

where:
\begin{itemize}[leftmargin=2em, itemsep=0pt]
  \item $\beta \in \R^p$ are regression coefficients and $\lambda$ is
        the skewness parameter;
  \item $z_{k, t}, \tau_{k, t}$ are knot-level latent processes
        governing skew and scale;
  \item $w_{k, t} \in \R^2$ is the (possibly time-varying) location of
        knot $k$ at time $t$;
  \item $g(s, t) = \arg\min_k \|s - w_{k, t}\|$ is the knot membership
        function (a Voronoi assignment);
  \item $C = C(\rho, \nu, \gamma)$ is the spatial covariance matrix on
        the observation grid and $L_C = \operatorname{chol}(C)^\top$.
\end{itemize}

We collect the time-$t$ slice of all latents into
$\Ltot_t = (z_{\cdot, t}, \tau_{\cdot, t}, w_{\cdot, t})$ and the full
fixed-parameter vector into $\theta$.

\subsection{AR(2) standardisation}

Each scalar latent process (one per knot per quantity $z, \tau, w$) is
generated by

\begin{equation}
  X_t = \phi_1 X_{t-1} + \phi_2 X_{t-2} + \eta_t,
  \qquad \eta_t \sim \Norm(0, \sigma_\eta^2),
  \label{eq:ar2}
\end{equation}

with $(\phi_1, \phi_2)$ in the stationarity region. We adopt the
construction in
\verb|simulate_ar2_standard()|, which fixes the
stationary variance and \emph{solves} for the innovation variance.

\begin{definition}[Standardised AR(2)]
\label{def:standardised}
Set $\gamma_0 = 1$, $\gamma_1 = \phi_1 / (1 - \phi_2)$, and
\begin{equation}
  \sigma_\eta^2
    = \gamma_0 - \phi_1 \gamma_1 - \phi_2 \gamma_2,
  \qquad \gamma_2 = \phi_1 \gamma_1 + \phi_2 \gamma_0.
  \label{eq:innovvar}
\end{equation}
The initial pair $(X_1, X_2)$ is drawn jointly from
\[
  \Norm\!\left(0,
    \begin{bmatrix} 1 & \gamma_1 \\ \gamma_1 & 1 \end{bmatrix}\right),
\]
and the recursion~\eqref{eq:ar2} runs forward for $t \ge 3$.
\end{definition}

\begin{assumption}[Stationarity]
\label{assn:stationary}
$(\phi_1, \phi_2)$ satisfies
$|\phi_2| < 1, \ \phi_1 + \phi_2 < 1, \ \phi_2 - \phi_1 < 1$,
so the roots of the AR(2) characteristic polynomial lie strictly
outside the unit circle.
\end{assumption}

\subsection{Prediction task}

The simstudy holdout fixes a subset $\mathcal S_{\mathrm{test}}$ of
test \emph{sites}; the training set is
$\D = \{Y_t(s) : s \notin \mathcal S_{\mathrm{test}}, \, t \in
\{1, \dots, T\}\}$, and the prediction targets are
$\{Y_t(s^*) : s^* \in \mathcal S_{\mathrm{test}}, \, t \in
\{1, \dots, T\}\}$. Note that every time point is observed at every
training site.

\begin{definition}[Brier score]
\label{def:brier}
For threshold $c$ and predictive probability
$\hat p(s^*, t; c) = \Pr(Y_t(s^*) > c \given \D)$,
\begin{equation}
  \mathrm{BS}(c) = \E\big[(\indicator\{Y_t(s^*) > c\}
                        - \hat p(s^*, t; c))^2\big].
  \label{eq:brier}
\end{equation}
\end{definition}

% ====================================================================
\section{Marginal Invariance of the Standardised AR(2)}
\label{sec:ar2-construction}

\begin{proposition}[Marginal invariance to $\phi$]
\label{prop:invariance}
Under Definition~\ref{def:standardised} and
Assumption~\ref{assn:stationary}, for every $t \in \{1, \dots, T\}$,
\begin{equation}
  X_t \sim \Norm(0, 1)
  \qquad \text{regardless of } (\phi_1, \phi_2).
  \label{eq:invariance}
\end{equation}
\end{proposition}

\begin{proof}
By construction $X_1, X_2 \sim \Norm(0, 1)$ marginally. Suppose
inductively that $\Var(X_{t-1}) = \Var(X_{t-2}) = 1$ and
$\Cov(X_{t-1}, X_{t-2}) = \gamma_1$. Then from~\eqref{eq:ar2},
\[
  \Var(X_t)
    = \phi_1^2 + \phi_2^2 + 2 \phi_1 \phi_2 \gamma_1 + \sigma_\eta^2.
\]
Substituting~\eqref{eq:innovvar} and using
$\gamma_2 = \phi_1 \gamma_1 + \phi_2$ gives
$\Var(X_t) = 1$. Mean is zero by linearity. Joint normality of $X_t$
with previous lags is preserved by the linear recursion.
\end{proof}

\begin{corollary}[AR(1) toy]
\label{cor:ar1}
For an AR(1) process standardised to $\gamma_0 = 1$,
$\sigma_\eta^2 = 1 - \phi^2$, and $X_t \sim \Norm(0, 1)$ for every
admissible $\phi$. The autocorrelation $\Cov(X_t, X_{t+h}) = \phi^{|h|}$
is the only quantity that moves with $\phi$.
\end{corollary}

\begin{remark}[Identifiability]
$\phi$ is identifiable only from \emph{joint} distributions across
time; the one-point marginal is $\phi$-blind. Sampling one $X_t$
alone provides zero information about $\phi$.
\end{remark}

% ====================================================================
\section{Why Brier Marginalises $\phi$ Away}
\label{sec:derivations}

\subsection{The predictive integral collapses to the time-$t$ marginal}

\begin{proposition}[Predictive at $(s^*, t)$ depends only on time-$t$
latents]
\label{prop:kernel-decomposition}
For the kernel~\eqref{eq:kernel}, the conditional CDF at $(s^*, t)$
given $(\theta, \Ltot_t)$ does not depend on $\Ltot_{t'}$ for any
$t' \ne t$. Consequently the marginal predictive is
\begin{equation}
  \hat p(s^*, t; c)
    = \int \!\!\int
        \Pr\!\big(Y_t(s^*) > c \given \theta, \Ltot_t\big) \,
        \pi(\Ltot_t \given \theta, \D) \,
        \pi(\theta \given \D) \, d\Ltot_t \, d\theta.
  \label{eq:predictive}
\end{equation}
\end{proposition}

\begin{proof}
Direct from inspection of~\eqref{eq:kernel}: the right-hand side
involves only $\Ltot_t$ at the same $t$. Fubini and the conditional
factorisation $\pi(\theta, \Ltot_t \given \D) = \pi(\Ltot_t \given
\theta, \D) \pi(\theta \given \D)$ give the displayed form.
\end{proof}

\begin{proposition}[$\phi$-free prior marginal at one $t$]
\label{prop:phi-free}
Under Definition~\ref{def:standardised}, the \emph{prior} marginal of
$\Ltot_t$ at any single $t$ does not depend on
$\phi := (\phi_z, \phi_\tau, \phi_w)$.
\end{proposition}

\begin{proof}
Each component of $\Ltot_t$ is obtained by a fixed monotone transform
(gamma inverse copula for $\tau$, half-normal inverse copula for $z$,
inverse probit for $w$) applied to a standardised AR(2) latent.
Proposition~\ref{prop:invariance} establishes that each underlying
latent has marginal $\Norm(0, 1)$ for every admissible $\phi$. The
fixed transforms preserve this invariance.
\end{proof}

\begin{corollary}[Bayes vs.\ prior]
\label{cor:posterior-route}
Combining Propositions~\ref{prop:kernel-decomposition} and
\ref{prop:phi-free}, the only route through which $\phi$ can affect
$\hat p(s^*, t; c)$ is the posterior $\pi(\theta \given \D)$ together
with the conditional $\pi(\Ltot_t \given \theta, \D)$. The
likelihood $p(\D \given \theta)$ couples observations across time and
therefore depends on $\phi$, but the kernel at the held-out cell does
not.
\end{corollary}

\subsection{Four leakage channels}

We enumerate the channels through which $\phi$ can move the predictive
at $(s^*, t)$ in spite of Corollary~\ref{cor:posterior-route}.

\begin{description}[leftmargin=2em, itemsep=0.25em]
  \item[Channel A: $z_{g, t}$ via temporal pooling.] AR(2) tightens
    $\pi(z_{k, t} \given \D)$ at training $t$ by borrowing from
    $z_{k, t \pm 1}, z_{k, t \pm 2}$. The benefit is bounded by how
    much information the time-$t$ training likelihood already
    provides.
  \item[Channel B: $\tau_{g, t}$ via temporal pooling.] Symmetric to
    Channel A.
  \item[Channel C: knot membership $g(s^*, t)$.] AR(2) tightens
    $\pi(w_{k, t} \given \D)$ at training $t$. Because $g(s^*, t)$ is a
    discrete \emph{argmin} over $\{w_{1, t}, \dots, w_{K, t}\}$, the
    posterior tightening propagates to $g$ only when the perturbation
    crosses a Voronoi boundary near $s^*$. For $K = 1$ the channel is
    closed entirely.
  \item[Channel D: spatial noise $L_C \varepsilon_t / \sqrt{\tau_{g,t}}$.]
    $\varepsilon_t$ is i.i.d.\ in $t$ given $\theta$. AR(2) does not
    appear in this channel at all.
\end{description}

We also account for a cost channel that operates against any benefit
from A--C.

\begin{description}[leftmargin=2em, itemsep=0.25em]
  \item[Channel E (cost): posterior on $\phi$.] Methods 7 and 8 turn
    on AR(2) for $z, \tau, w$ simultaneously, introducing six new
    parameters $\phi_{1z}, \phi_{2z}, \phi_{1\tau}, \phi_{2\tau},
    \phi_{1w}, \phi_{2w}$. Their joint posterior contributes
    variability to the predictive integral~\eqref{eq:predictive}.
\end{description}

\subsection{Variance decomposition}

\begin{lemma}[Predictive variance decomposition]
\label{lem:variance}
For the held-out predictive probability
$\hat F_{s^*, t}(c) = \hat p(s^*, t; c)$,
\begin{equation}
  \Var\!\big(\hat F_{s^*, t}(c)\big)
    \approx \E_\phi\!\big[\Var\!\big(\hat F \given \phi\big)\big]
          + \Var_\phi\!\big(\E[\hat F \given \phi]\big).
  \label{eq:variance-decomp}
\end{equation}
\end{lemma}

The AR(2) prior reduces the first summand through Channels A--C; it
introduces the second summand, which is identically zero under
i.i.d.\ time. The Brier score is sensitive to the sum, not the
decomposition. Empirically (Section~\ref{sec:evidence}) the second
summand dominates the first reduction in the simstudy configuration.

\subsection{Evidence in the simstudy}
\label{sec:evidence}

Table~\ref{tab:rel-brier} summarises Brier scores relative to the
Gaussian baseline (method 1) at selected quantiles for setting 9.

\begin{table}[h]
\centering
\caption{Relative Brier score versus method 1, setting 9. Smaller is
         better; entries below~1 indicate improvement over the
         Gaussian baseline.}
\label{tab:rel-brier}
\begin{tabular}{c r r r r}
\toprule
$q$ & M2 (skew-$t$, $K{=}1$) & M4 (skew-$t$, $K{=}5$)
    & M7 (AR(2), $K{=}1$)    & M8 (AR(2), $K{=}5$) \\
\midrule
0.90  & 0.961 & 0.973 & 0.968 & 0.972 \\
0.95  & 0.951 & 0.969 & 0.964 & 0.969 \\
0.97  & 0.943 & 0.965 & 0.958 & 0.958 \\
0.99  & 0.955 & 0.987 & 0.975 & 0.980 \\
0.995 & 0.946 & 0.955 & 0.957 & 0.962 \\
\bottomrule
\end{tabular}
\end{table}

The i.i.d.\ skew-$t$ at $K = 1$ (method 2) systematically beats its
AR(2) sibling (method 7). Setting~10 shows an even sharper version of
the same phenomenon: method 8 ($K = 5$ with AR(2)) ranges from
\(0.96\) to \(1.16\) of the Gaussian baseline, while method 4
($K = 5$ without AR(2)) is at \(0.87\)--\(0.94\).

% ====================================================================
\section{Redesigns That Restore the AR(2) Signal}
\label{sec:redesign}

Three designs expose the temporal coupling that the spatial-only Brier
ignores.

\subsection{Design A: block-time-forecast holdout}

Hold out the last $T_h$ time points; train on
$\{Y_{t}(s) : t \le T - T_h, \, s \in \mathcal S\}$. The predictive at
horizon $t^* > T - T_h$ requires the conditional
\[
  z_{k, t^*} \given \D \;\sim\;
    \Norm\!\big(\phi_1 \hat z_{k, t^* - 1} + \phi_2 \hat z_{k, t^* - 2},\,
                \sigma_\eta^2 + \text{filter uncertainty}\big),
\]
which is strictly $\phi$-dependent. Brier and any scoring rule applied
to the forecast cells will distinguish AR(2) from i.i.d.\ time.

\subsection{Design B: multivariate score on the spatial holdout}

Keep the spatial holdout. Score the held-out site-level trajectory
$\bm y_{s^*} = (Y_1(s^*), \dots, Y_T(s^*))$ as a $T$-vector.

\paragraph{Energy score.} For predictive $F$ and observation
$\bm y \in \R^T$,
\begin{equation}
  \mathrm{ES}(F, \bm y)
    = \E_F \|X - \bm y\|_2
      - \tfrac{1}{2} \E_F \|X - X'\|_2,
  \qquad X, X' \overset{\text{iid}}{\sim} F.
\end{equation}
The Monte Carlo estimator with $M$ samples is
\[
  \widehat{\mathrm{ES}}
    = \frac{1}{M} \sum_{m=1}^{M} \|X^{(m)} - \bm y\|_2
      - \frac{1}{2 M (M - 1)} \sum_{m \ne m'} \|X^{(m)} - X^{(m')}\|_2.
\]

\paragraph{Variogram score of order $p$.}
\begin{equation}
  \mathrm{VS}_p(F, \bm y)
    = \sum_{t < t'} w_{t, t'}
      \Big( |y_t - y_{t'}|^p - \E_F |X_t - X_{t'}|^p \Big)^2.
\end{equation}
With $p = 0.5$ the score is particularly sensitive to temporal
dependence misspecification.

\subsection{Design C (recommended): combined holdout + energy score}

Hold out both a set of test sites and the last $T_h$ time points at
the remaining sites. Report:
\begin{enumerate}[leftmargin=2em, itemsep=0pt]
  \item marginal Brier on $(s^*, t)$ for $s^* \in
        \mathcal S_{\mathrm{test}}$, $t \le T - T_h$ --- the legacy
        comparison;
  \item energy score on $(s, t_{\mathrm{block}})$ for $s$ in the
        training spatial set, $t > T - T_h$ --- the temporal-coupling
        test;
  \item energy score on $(s^*, t_{\mathrm{block}})$ ---
        space-and-time extrapolation.
\end{enumerate}

\begin{remark}[Reporting]
Report paired score differences across the 50 simulation replicates
and bootstrap the resulting boxplots. Also report
score-per-CPU-hour: methods 7 and 8 are more expensive, and a small
score gain may not be worth the additional compute.
\end{remark}

% ====================================================================
\section{Intuition Summary}
\label{sec:intuition}

AR(2) is a model of how today's latent values relate to yesterday's
and the day before yesterday's. The Brier score on the spatial
holdout asks for the marginal exceedance probability at a new
\emph{location} on a day that is already observed at other locations.
The temporal coupling that AR(2) models has no opportunity to act,
because every relevant time point is already pinned by training-time
likelihood. The score sees only the cost of the additional
parameters --- six $\phi$'s with their own posterior variance --- and
penalises the model for buying capability it never gets to use.

The fix is not to remove AR(2); it is to design an experiment that
asks the model to act like a temporal model. Forecast forward in
time, or score whole trajectories at once. The temporal coupling then
enters the predictive integral non-trivially, and the AR(2) prior
recovers its rent.

% ====================================================================
\section{Conclusion and Future Work}
\label{sec:conclusion}

We have established that the AR(2) standardisation
(Definition~\ref{def:standardised}) makes the prior time-$t$ marginal
of every latent invariant to the autoregressive coefficients
(Proposition~\ref{prop:invariance}, Proposition~\ref{prop:phi-free}).
The spatial-only Brier score integrates the predictive kernel
against precisely this marginal
(Proposition~\ref{prop:kernel-decomposition}), so the temporal
coupling vanishes from the loss. Of the four leakage channels through
which $\phi$ could re-enter via the posterior, three (Channels A--C)
are damped by likelihood dominance, discrete knot membership, or
spatial $\phi$-independence; the fourth (Channel D) is closed. The
cost channel (E) is unbounded above and inflates the predictive
variance, which Brier penalises directly. The net effect is that
i.i.d.-time skew-$t$ wins the comparison on this score.

Future work, in order of priority:
\begin{enumerate}[leftmargin=2em, itemsep=0pt]
  \item Implement Design C in
        \texttt{scores.R} and \texttt{run-settings.R}: extend the MCMC
        predictive loop in \texttt{rpotspatTS\_arp()} to propagate the
        AR(2) recursion past the training horizon, and add
        \texttt{EnergyScore} and \texttt{VariogramScore} alongside
        \texttt{BrierScore}.
  \item Re-run settings 9 and 10 with $T_h \in \{3, 5, 10\}$ and
        report paired score differences between methods (2, 7) and
        (4, 8).
  \item Run model selection over AR order $p \in \{0, 1, 2\}$ using
        DIC/WAIC before fixing the temporal-prior order, so any
        observed AR(2) advantage cannot be attributed to a
        coincidentally matched simulation truth.
  \item Investigate whether channel C activates non-trivially when
        the data-generating $w$ has small inter-knot separation,
        i.e.\ when many held-out sites sit near Voronoi boundaries.
\end{enumerate}

% ====================================================================
\appendix

\section{Reference R implementation}
\label{app:r-code}

The standardised AR(2) simulator in
\path{code/00_core/auxfunctions_ar2.R}.

\begin{lstlisting}
simulate_ar2_standard <- function(nt, n, phi1, phi2,
                                  tol = 1e-06, param_name = NULL) {
  # ... stationarity check ...
  denom  <- 1 - phi2
  gamma0 <- 1
  gamma1 <- phi1 * gamma0 / denom
  gamma2 <- phi1 * gamma1 + phi2 * gamma0
  sigma2 <- gamma0 - phi1 * gamma1 - phi2 * gamma2
  innov.sd <- sqrt(sigma2)

  ar2 <- matrix(0, nrow = n, ncol = nt)
  init.states <- draw_ar2_initial(n, gamma0, gamma1)
  ar2[, 1] <- init.states[, 1]
  ar2[, 2] <- init.states[, 2]

  for (t in 3:nt) {
    ar2[, t] <- phi1 * ar2[, t - 1] +
                phi2 * ar2[, t - 2] +
                rnorm(n, 0, innov.sd)
  }
  ar2
}
\end{lstlisting}

The construction matches Definition~\ref{def:standardised}: the
marginal variance is fixed at $\gamma_0 = 1$ and the innovation
variance is derived. Initial states are drawn jointly from the
stationary bivariate Gaussian via \texttt{draw\_ar2\_initial()},
removing the need for a burn-in.

\section{Verification by reverse simulation}
\label{app:verification}

To confirm the marginal invariance empirically we simulate $n = 1000$
AR(2) paths under
$(\phi_1, \phi_2) = (0.6, -0.3)$ via manual recursion and via
\texttt{stats::filter(\ldots, method = "recursive")}, fit
\texttt{arima(order = c(2, 0, 0))} to each, and run a residual
diagnostic. Both fits recover the true coefficients within one
standard error, and the Ljung--Box and Shapiro--Wilk tests on the
residuals are non-significant ($p > 0.4$). Manual recursion and
filtering produce statistically equivalent AR(2) series.

\end{document}
